\documentclass[11pt,a4paper]{article}
\usepackage[utf8]{inputenc}
\usepackage[T1]{fontenc}
\usepackage{amsmath,amssymb,amsfonts}
\usepackage{booktabs}
\usepackage{hyperref}
\usepackage{geometry}
\usepackage{enumitem}
\usepackage{algorithm}
\usepackage{algorithmic}
\usepackage{graphicx}
\usepackage{xcolor}
\usepackage{tabularx}
\usepackage{amsthm}
\newtheorem{definition}{Definition}

\geometry{margin=2.5cm}

\title{\textbf{RDD-Video: Resolution-Drift Diffusion for Spatiotemporal Multi-Resolution Video Generation}\\[0.5em]\large Technical Proposal}
\author{Research Notes}
\date{\today}

\begin{document}
\maketitle

\begin{abstract}
We propose RDD-Video, an extension of the Resolution-Drift Diffusion (RDD) framework from 2D image generation to video generation. The core innovation is a \textbf{spatiotemporal drift} that simultaneously encodes spatial resolution progression and temporal frame-rate progression within a single continuous Markov process. Unlike existing multi-resolution video methods (Pyramid-Flow, TPDiff) that require explicit renoising at stage boundaries, RDD-Video maintains mathematical continuity via its drift mechanism, enabling compatibility with high-order ODE solvers (DPM-Solver) and achieving $10$--$20\times$ fewer function evaluations than SDE-based methods. This document provides a complete technical proposal including architecture design, mathematical formulation, training strategy, sampling algorithm, and experiment plan.
\end{abstract}

\tableofcontents
\newpage

%=============================================================================
\section{Related Work: Multi-Resolution Video Generation}
%=============================================================================

\subsection{Pyramid-Flow (ICLR 2025)}

\textbf{Core idea:} Divides the Flow Matching trajectory into 3 spatial pyramid stages operating at $\frac{1}{4}\times$, $\frac{1}{2}\times$, and $1\times$ spatial resolution. Each stage uses rectified flow with start/end points defined by upsampled latents + noise. Stage transitions use \emph{renoising}: adding block noise to match the distribution of the next stage's starting point.

\textbf{Key formulas:}
\begin{align}
x_t &= \text{ratio} \cdot \text{start\_point} + (1-\text{ratio}) \cdot \text{end\_point} \\
\text{target} &= \text{start\_point} - \text{end\_point} \quad \text{(rectified flow velocity)}
\end{align}
Stage transition (renoising):
\begin{equation}
\text{latents}_\text{new} = \alpha \cdot \text{Up}(\text{latents}_\text{old}) + \beta \cdot \text{block\_noise}
\end{equation}
where $\alpha, \beta$ are functions of a hyperparameter $\gamma$ controlling noise strength.

\textbf{Architecture:} Single DiT (FLUX-based), 768p video at 24 FPS up to 10s. Uses temporal AR with causal attention.

\textbf{Limitations:}
\begin{itemize}[nosep]
\item Explicit renoising at stage boundaries causes information loss
\item Each stage is an independent flow $\to$ no global mathematical unification
\item Incompatible with DPM-Solver (cannot cross stage boundaries with high-order steps)
\item AR temporal generation limits inference speed
\end{itemize}

\textbf{Code:} \texttt{github.com/jy0205/Pyramid-Flow} --- key file: \texttt{pyramid\_dit\_for\_video\_gen\_pipeline.py}. The \texttt{add\_pyramid\_noise()} function (L300--349) implements the multi-stage flow matching training: for each stage $i_s$, it constructs start/end points via spatial upsampling + noise mixing, then linearly interpolates to get $x_t$.

\subsection{TPDiff --- Temporal Pyramid Diffusion (ICLR 2026)}

\textbf{Core idea:} Progressively increases \emph{frame rate} (temporal resolution) along the diffusion process. Early high-noise stages use $\frac{1}{4}$ frame rate; only the final stage operates at full frame rate. Uses \emph{stage-wise diffusion} training with data-noise alignment.

\textbf{Key formulas:}
Stage-wise endpoints:
\begin{align}
\hat{x}_{s_k} &= \gamma_{s_k} \text{Up}(\text{Down}(x_0, 2^{k+1}), 2) + \sigma_{s_k} \epsilon \\
\hat{x}_{e_k} &= \gamma_{e_k} \text{Down}(x_0, 2^k) + \sigma_{e_k} \epsilon
\end{align}
where Down/Up are \emph{temporal} downsampling/upsampling.

Stage-wise noise (derived from ODE decomposition):
\begin{equation}
\epsilon_k = \frac{\hat{x}_{e_k}/\gamma_{e_k} - \hat{x}_{s_k}/\gamma_{s_k}}{\sigma_{e_k}/\gamma_{e_k} - \sigma_{s_k}/\gamma_{s_k}}
\end{equation}

Inference renoising:
\begin{equation}
\hat{x}_{s_{k-1}} = \frac{\sqrt{2}\gamma_{s_k}}{\sigma_{s_k} + \sqrt{2}\gamma_{s_k}} \text{Up}(\hat{x}_{e_k}) + \frac{\sqrt{2}\sigma_t}{2} n', \quad n' \sim \mathcal{N}(0, \Sigma')
\end{equation}
with anti-correlated noise $\Sigma' = \begin{pmatrix}1 & -1 \\ -1 & 1\end{pmatrix}$ (block structure).

\textbf{Key innovation:} Data-noise alignment via Hungarian algorithm minimizes $\text{dist}(x, \epsilon)$ within a batch, making the ODE path approximately deterministic $\to$ 2$\times$ faster training.

\textbf{Results:} 50\% training cost reduction, 1.5$\times$ inference speedup on both DDIM and Flow Matching.

\textbf{Limitations:}
\begin{itemize}[nosep]
\item Still requires renoising between temporal stages
\item Only addresses temporal dimension; spatial is fixed
\item Incompatible with DPM-Solver across stage boundaries
\end{itemize}

\textbf{Code:} Not publicly available. Method described in arXiv:2503.09566.

\subsection{Open-Sora Plan (2025)}

\textbf{Core components:}
\begin{itemize}[nosep]
\item \textbf{WF-VAE}: Wavelet-Flow VAE with 4$\times$8$\times$8 compression (temporal$\times$spatial)
\item \textbf{Skiparse Denoiser}: 3D DiT + Skiparse Attention + 3D RoPE, 2.7B params
\item \textbf{Progressive training}: Image $\to$ low-res video $\to$ high-res video
\end{itemize}

\textbf{Skiparse Attention:} Reduces $O(N^2)$ to $O(N^2/k)$ via sparse token grouping (sparse ratio $k=4$). Alternates between Single Skip (every $k$-th token) and Group Skip (contiguous groups).

\textbf{Relevance to RDD-Video:} Skiparse can serve as the backbone attention mechanism for RDD-Video's DiT, naturally supporting variable sequence lengths across resolution stages.

\subsection{CogVideoX (2024)}

Uses \textbf{progressive training} (low-res $\to$ high-res) and 3D Full Attention with Expert AdaLN. Key insight: joint image-video training with 3D Causal VAE (4$\times$8$\times$8 compression).

\subsection{Matryoshka Diffusion Models (ICLR 2024)}

Joint multi-resolution diffusion in extended space with NestedUNet. Demonstrates text-to-video on WebVid at 16 frames. Uses resolution-dependent shifted cosine noise schedule.

\subsection{Jenga / ProRes (2025)}

Training-free progressive resolution strategy for inference: starts denoising at low spatial resolution, upsamples as denoising proceeds. Achieves 8.83$\times$ speedup on VBench with only 0.01\% quality drop.

%=============================================================================
\section{RDD 2D: Core Method Recap}
%=============================================================================

\subsection{Forward Process with Drift}
\begin{equation}
q(x_t | x_0) = \mathcal{N}\left(\sqrt{\bar\alpha_t} \cdot x_0 + \sqrt{\bar\alpha_t} \cdot \kappa(t) \cdot d, \; (1-\bar\alpha_t)\mathbf{I}\right)
\end{equation}
where:
\begin{itemize}[nosep]
\item $d = \text{Up}(\text{Down}(x_0)) - x_0$ is the drift vector (negative high-frequency residual)
\item $\kappa(t) \in [0,1]$ is a piecewise linear schedule within each stage
\item Stages defined by \texttt{stage\_lengths}, e.g., $[300, 300, 400]$ for total 1000 steps
\end{itemize}

\subsection{Key Properties}
\begin{enumerate}[nosep]
\item \textbf{Single Markov chain}: No stage boundaries, drift varies continuously
\item \textbf{Closed-form posterior}: $q(x_{t-1}|x_t, x_0)$ is Gaussian with drift correction
\item \textbf{Single prediction target}: $\tilde\varepsilon = \varepsilon + \sqrt{\bar\alpha_t/(1-\bar\alpha_t)} \cdot \kappa(t) \cdot d$
\item \textbf{DPM-Solver compatible}: $x_{t-1} = \text{dpm\_step}(x_t) + \sqrt{\bar\alpha_{t-1}} \cdot (\kappa(t-1) - \kappa(t)) \cdot d$
\item \textbf{Frequency-adaptive SNR}: Low-freq SNR unchanged, high-freq SNR decays by $(1-\kappa(t))^2$
\end{enumerate}

\subsection{Sampling (2D)}
\begin{enumerate}[nosep]
\item Stage 1 (8$\times$8): DDPM/DPM-Solver from pure noise to 8$\times$8 latent
\item Stage transition: $x_\text{next} = \sqrt{\bar\alpha_{t_\text{start}}} \cdot \text{Up}(\hat{x}_0) + \sqrt{1-\bar\alpha_{t_\text{start}}} \cdot \varepsilon$ (noise filling)
\item Stage 2 (16$\times$16): Continue denoising from $t_\text{start}$
\item Stage transition + Stage 3 (32$\times$32): Continue to completion
\end{enumerate}

%=============================================================================
\section{RDD-Video: Spatiotemporal Drift Design}
%=============================================================================

\subsection{Core Principle: Spatiotemporal Drift}

The key insight is that video has two dimensions of resolution: \textbf{spatial} (pixel resolution) and \textbf{temporal} (frame rate). RDD's drift mechanism naturally extends to both:

\begin{definition}[Spatiotemporal Drift]
For a video latent $x_0 \in \mathbb{R}^{C \times T \times H \times W}$, define the spatiotemporal drift:
\begin{equation}
d = \text{STUp}(\text{STDown}(x_0)) - x_0
\end{equation}
where:
\begin{align}
\text{STDown}(x_0) &= \text{TemporalDown}_{\tau}(\text{SpatialDown}_s(x_0)) \in \mathbb{R}^{C \times T/\tau \times H/s \times W/s} \\
\text{STUp}(\cdot) &= \text{TemporalUp}_{\tau}(\text{SpatialUp}_s(\cdot)) \in \mathbb{R}^{C \times T \times H \times W}
\end{align}
with $s$ the spatial downsampling factor and $\tau$ the temporal downsampling factor.
\end{definition}

This drift captures \emph{both} spatial high-frequency detail and temporal high-frequency motion simultaneously.

\subsection{Multi-Stage Spatiotemporal Resolution Pyramid}

We define $K$ stages with progressively increasing spatiotemporal resolution:

\begin{table}[h]
\centering
\begin{tabular}{cccccc}
\toprule
Stage $k$ & Spatial Res & Temporal Res & Spatial Factor & Temporal Factor & Timestep Range \\
\midrule
1 (coarsest) & $H/4 \times W/4$ & $T/4$ frames & $s=4$ & $\tau=4$ & $[0, T_1)$ \\
2 & $H/2 \times W/2$ & $T/2$ frames & $s=2$ & $\tau=2$ & $[T_1, T_2)$ \\
3 (finest) & $H \times W$ & $T$ frames & $s=1$ & $\tau=1$ & $[T_2, T_3)$ \\
\bottomrule
\end{tabular}
\caption{3-stage spatiotemporal pyramid. Example: \texttt{stage\_lengths} = [250, 350, 400] for 1000 total steps.}
\end{table}

\subsection{Hierarchical Drift Definition}

For the $K$-stage pyramid, we define a \textbf{cumulative drift} at each stage boundary:

\textbf{Stage $k$ drift vector:}
\begin{equation}
d_k = \text{STUp}_{k \to k+1}(\text{STDown}_{k+1 \to k}(x_0^{(k+1)})) - x_0^{(k+1)}
\end{equation}
where $x_0^{(k+1)}$ is the full-resolution clean latent at stage $k+1$'s resolution.

\textbf{Unified single-chain forward process:}
\begin{equation}
q(x_t | x_0) = \mathcal{N}\left(\sqrt{\bar\alpha_t} \cdot x_0^{(k)} + \sqrt{\bar\alpha_t} \cdot \kappa_k(t) \cdot d_k, \; (1-\bar\alpha_t)\mathbf{I}\right)
\end{equation}
where $k$ is the stage active at timestep $t$ (determined by stage boundaries), and $x_0^{(k)}$ is the target latent at stage $k$'s resolution.

\subsection{Temporal Drift Specifics}

\textbf{Temporal downsampling options:}
\begin{itemize}[nosep]
\item Average pooling along temporal axis (analogous to spatial avg-pool)
\item Stride-$\tau$ subsampling (select every $\tau$-th frame)
\item Temporal low-pass filter (Gaussian blur along time) --- \emph{recommended}
\end{itemize}

\textbf{Temporal upsampling for noise filling:}
\begin{itemize}[nosep]
\item Nearest-neighbor repetition (simplest, used by TPDiff)
\item Linear interpolation between adjacent frames
\item Slerp interpolation in latent space
\end{itemize}

The temporal drift $d_\text{temporal} = \text{TemporalUp}(\text{TemporalDown}(x_0)) - x_0$ captures inter-frame high-frequency motion (fast movements, temporal edges).

\subsection{Mathematical Properties Preserved}

\begin{enumerate}[nosep]
\item \textbf{Single Markov chain}: The spatiotemporal drift is still a deterministic function of $x_0$, so the forward process remains a valid Gaussian perturbation kernel.
\item \textbf{Closed-form posterior}: Since $d$ is fixed given $x_0$, the posterior $q(x_{t-1}|x_t,x_0)$ remains tractable:
\begin{equation}
\text{posterior\_mean} = \text{standard\_posterior\_mean} + \text{coef\_drift}(t) \cdot d
\end{equation}
\item \textbf{DPM-Solver compatibility}: The correction step becomes:
\begin{equation}
x_{t-1} = \text{dpm\_step}(x_t) + \sqrt{\bar\alpha_{t-1}} \cdot (\kappa(t-1) - \kappa(t)) \cdot d
\end{equation}
This works identically in the video setting since $d$ is a 3D tensor.
\item \textbf{Frequency-adaptive SNR (spatiotemporal)}:
\begin{equation}
\text{SNR}_\text{RDD}(f_s, f_t, t) = \frac{\bar\alpha_t \cdot [(1-\kappa(t)) + \kappa(t) \cdot \text{LP}(f_s, f_t)]^2}{1-\bar\alpha_t}
\end{equation}
where $\text{LP}(f_s, f_t)$ is the spatiotemporal low-pass filter response. High spatial \emph{and} high temporal frequencies get additional SNR decay.
\end{enumerate}

%=============================================================================
\section{Architecture Design}
%=============================================================================

\subsection{Overall Architecture}

\begin{center}
\fbox{\parbox{0.9\textwidth}{
\textbf{RDD-Video Architecture:}\\[0.3em]
Video $\xrightarrow{\text{3D VAE}}$ Latent $z_0 \in \mathbb{R}^{C \times T' \times H' \times W'}$ $\xrightarrow{\text{RDD Forward + ST Drift}}$ Noisy latent $z_t$ $\xrightarrow{\text{3D DiT}}$ Predicted $\tilde\varepsilon$ $\xrightarrow{\text{RDD Reverse}}$ Clean latent $\hat{z}_0$ $\xrightarrow{\text{3D VAE Decoder}}$ Video
}}
\end{center}

\subsection{3D VAE (Encoder/Decoder)}

Use an existing high-quality video VAE:
\begin{itemize}[nosep]
\item \textbf{Option A}: WF-VAE (Open-Sora Plan) --- compression $4\times8\times8$, wavelet-based
\item \textbf{Option B}: CogVideoX 3D Causal VAE --- compression $4\times8\times8$
\item \textbf{Option C}: SD3-based VAE + temporal extension (Pyramid-Flow approach)
\end{itemize}

All options yield latents of shape $C \times T/4 \times H/8 \times W/8$ (e.g., 24 frames $\to$ 6 temporal tokens for a 2-second clip at 12 FPS).

\subsection{3D DiT Backbone}

\textbf{Unified model} approach (single DiT handles all resolutions):

\begin{itemize}[nosep]
\item \textbf{Base}: DiT-XL/2 architecture extended to 3D
\item \textbf{Attention}: Skiparse Attention (from Open-Sora Plan) with sparse ratio $k=4$
  \begin{itemize}[nosep]
  \item First/last 2 layers: Full 3D attention
  \item Middle layers: Skiparse (alternating Single Skip and Group Skip)
  \end{itemize}
\item \textbf{Position encoding}: 3D RoPE (Rotary Position Embedding)
  \begin{itemize}[nosep]
  \item Supports variable sequence lengths across stages
  \item Each spatial/temporal position gets independent frequency encoding
  \end{itemize}
\item \textbf{Conditioning}: AdaLN-Zero with class/text embedding
\item \textbf{Patch embedding}: $2\times2\times1$ spatial patches (temporal is handled by VAE compression)
\end{itemize}

\textbf{Key design choice --- Unified vs. Multistage:}
Based on 2D RDD experiments showing Unified achieves the best FID, we adopt the \textbf{unified} approach: a single 3D DiT-XL handles all spatiotemporal resolutions. The 3D RoPE naturally adapts to different token counts per stage.

\subsection{Model Size Estimates}

\begin{table}[h]
\centering
\begin{tabular}{lcc}
\toprule
Component & Parameters & Notes \\
\midrule
3D VAE & $\sim$200M & Frozen (pretrained) \\
3D DiT-XL & $\sim$700M & Trainable \\
Text Encoder (T5-XXL) & $\sim$4.7B & Frozen \\
\midrule
Total trainable & $\sim$700M & \\
\bottomrule
\end{tabular}
\end{table}

%=============================================================================
\section{Training Strategy}
%=============================================================================

\subsection{Multi-Stage Training Objective}

At each training step, for each stage $k$:
\begin{enumerate}[nosep]
\item Obtain full-resolution video latent $z_0$ from VAE
\item Compute spatiotemporal downsampled latent $z_0^{(k)}$ at stage $k$'s resolution
\item Compute drift $d_k$ at stage $k$
\item Sample $t \sim \text{Uniform}[T_k, T_{k+1})$
\item Compute $\kappa_k(t) = (t - T_k)/(T_{k+1} - T_k)$
\item Add noise with drift: $z_t = \sqrt{\bar\alpha_t} (z_0^{(k)} + \kappa_k(t) \cdot d_k) + \sqrt{1-\bar\alpha_t} \cdot \varepsilon$
\item Predict $\tilde\varepsilon_\theta(z_t, t)$ with the 3D DiT
\item Compute loss: $\|\tilde\varepsilon_\theta - (\varepsilon + \sqrt{\bar\alpha_t/(1-\bar\alpha_t)} \cdot \kappa_k(t) \cdot d_k)\|^2$
\end{enumerate}

\textbf{Total loss:} Sum over all stages with equal weight (same as 2D RDD Unified).

\subsection{Progressive Training Schedule}

Following CogVideoX and Open-Sora Plan's success:
\begin{enumerate}[nosep]
\item \textbf{Phase 1} (Image pretraining): Train 2D RDD on images (single frame, spatial-only drift). Initialize from DiT-XL ImageNet checkpoint.
\item \textbf{Phase 2} (Low-res video): Extend to short, low-resolution videos (16 frames, 256$\times$256). Add temporal layers and temporal drift.
\item \textbf{Phase 3} (High-res video): Scale to full resolution (64+ frames, 512$\times$512 or 768p). Activate all spatiotemporal stages.
\end{enumerate}

\subsection{Data-Noise Alignment (from TPDiff)}

Incorporate TPDiff's key innovation: use Hungarian algorithm to align noise samples with training data within each mini-batch:
\begin{algorithmic}
\STATE \texttt{assign\_mat} $\leftarrow$ \texttt{scipy.optimize.linear\_sum\_assignment(dist(x, $\epsilon$))}
\STATE $\epsilon' \leftarrow \epsilon$[\texttt{assign\_mat}]
\end{algorithmic}
This reduces variance in the ODE path estimation, accelerating convergence by $\sim$2$\times$.

\subsection{Training Hyperparameters}

\begin{table}[h]
\centering
\begin{tabular}{ll}
\toprule
Hyperparameter & Value \\
\midrule
Optimizer & AdamW ($\beta_1=0.9, \beta_2=0.999$) \\
Learning rate & $1\times10^{-4}$ \\
Batch size (global) & 256 (via gradient accumulation) \\
EMA decay & 0.9999 \\
Diffusion steps & 1000 \\
Stage lengths (3 stages) & [250, 350, 400] \\
Resolutions (spatial latent) & [8, 16, 32] (in VAE latent space) \\
Temporal resolutions & [$T/4$, $T/2$, $T$] frames \\
\bottomrule
\end{tabular}
\end{table}

%=============================================================================
\section{Sampling Strategy}
%=============================================================================

\subsection{RDD-Video Sampling Algorithm}

\begin{algorithm}[h]
\caption{RDD-Video Sampling (3 stages)}
\begin{algorithmic}[1]
\STATE \textbf{Input:} Text prompt $c$, CFG scale $w$, stage boundaries $[T_1, T_2, T_3]$
\STATE \textbf{Output:} Generated video $v$
\STATE
\STATE \textcolor{blue}{// Stage 1: Coarsest spatiotemporal resolution}
\STATE $z_{T_1} \sim \mathcal{N}(0, I)$ with shape $C \times T/4 \times H/4 \times W/4$
\STATE $\hat{z}_0^{(1)} \leftarrow \text{DPM-Solver}(z_{T_1}, \text{steps}=T_1, \text{model}=\theta)$
\STATE
\STATE \textcolor{blue}{// Stage transition 1$\to$2: Spatiotemporal noise filling}
\STATE $z_0^\text{up} \leftarrow \text{STUp}_{1\to2}(\hat{z}_0^{(1)})$ \quad // Upsample to $C \times T/2 \times H/2 \times W/2$
\STATE $z_{T_2-T_1} \leftarrow \sqrt{\bar\alpha_{t_\text{start}}} \cdot z_0^\text{up} + \sqrt{1-\bar\alpha_{t_\text{start}}} \cdot \varepsilon$ \quad // Noise filling
\STATE
\STATE \textcolor{blue}{// Stage 2: Mid resolution}
\STATE $\hat{z}_0^{(2)} \leftarrow \text{DPM-Solver}(z_{T_2-T_1}, \text{steps from } T_2-T_1 \text{ to } 0, \text{model}=\theta)$
\STATE
\STATE \textcolor{blue}{// Stage transition 2$\to$3: Spatiotemporal noise filling}
\STATE $z_0^\text{up} \leftarrow \text{STUp}_{2\to3}(\hat{z}_0^{(2)})$ \quad // Upsample to $C \times T \times H \times W$
\STATE $z_{T_3-T_2} \leftarrow \sqrt{\bar\alpha_{t_\text{start}}} \cdot z_0^\text{up} + \sqrt{1-\bar\alpha_{t_\text{start}}} \cdot \varepsilon$ \quad // Noise filling
\STATE
\STATE \textcolor{blue}{// Stage 3: Full resolution}
\STATE $\hat{z}_0^{(3)} \leftarrow \text{DPM-Solver}(z_{T_3-T_2}, \text{steps from } T_3-T_2 \text{ to } 0, \text{model}=\theta)$
\STATE
\STATE $v \leftarrow \text{VAE.decode}(\hat{z}_0^{(3)})$
\STATE \textbf{return} $v$
\end{algorithmic}
\end{algorithm}

\subsection{DPM-Solver with Spatiotemporal Drift Correction}

At each DPM-Solver step within a stage, apply the drift correction:
\begin{equation}
z_{t-1} = \text{dpm\_step}(z_t, \tilde\varepsilon_\theta) + \sqrt{\bar\alpha_{t-1}} \cdot (\kappa(t-1) - \kappa(t)) \cdot d
\end{equation}

During sampling, $d$ is estimated from the current prediction:
\begin{equation}
\hat{d} = \text{STUp}(\text{STDown}(\hat{z}_0)) - \hat{z}_0, \quad \hat{z}_0 = \frac{z_t - \sqrt{1-\bar\alpha_t} \cdot \tilde\varepsilon_\theta}{\sqrt{\bar\alpha_t}} - \kappa(t) \cdot \hat{d}
\end{equation}
This is solved iteratively (converges in 1--2 iterations in practice for 2D RDD).

\subsection{Noise Filling at Stage Transitions}

The ``constant upsample'' (noise filling) strategy that proved best in 2D RDD:
\begin{equation}
z_\text{next} = \sqrt{\bar\alpha_{t_\text{start}}} \cdot \text{STUp}(\hat{z}_0^{(k)}) + \sqrt{1-\bar\alpha_{t_\text{start}}} \cdot \varepsilon
\end{equation}
where $t_\text{start}$ is determined by the stage boundary. This is equivalent to treating the upsampled result as if it were the clean signal and forward-diffusing to the appropriate noise level.

\subsection{Classifier-Free Guidance}

Stage-dependent CFG (inspired by FreCaS):
\begin{itemize}[nosep]
\item Stage 1 (coarse structure): $w_1 = 7.0$ (strong guidance for global composition)
\item Stage 2 (mid detail): $w_2 = 4.0$
\item Stage 3 (fine detail): $w_3 = 2.0$ (weak guidance to preserve diversity/detail)
\end{itemize}

\subsection{Step Allocation}

Total steps budget: 80 (with DPM-Solver-2):
\begin{itemize}[nosep]
\item Stage 1: 20 steps (coarse, small tensor)
\item Stage 2: 30 steps (medium)
\item Stage 3: 30 steps (fine, largest tensor)
\end{itemize}

Computation is dominated by Stage 3 (full resolution), but RDD's drift ensures Stage 3 starts from a high-quality low-frequency estimate, requiring fewer refinement steps.

%=============================================================================
\section{Advantages over Existing Methods}
%=============================================================================

\begin{table}[h]
\centering
\small
\begin{tabularx}{\textwidth}{l|X|X|X|X}
\toprule
Property & \textbf{RDD-Video} & Pyramid-Flow & TPDiff & Cascade \\
\midrule
Single Markov chain & \checkmark & $\times$ & $\times$ & $\times$ \\
DPM-Solver compatible & \checkmark & $\times$ & $\times$ & $\times$ \\
No renoising needed & \checkmark & $\times$ & $\times$ & N/A \\
Single prediction target & \checkmark & \checkmark & \checkmark & $\times$ \\
Joint spatial+temporal pyramid & \checkmark & Partial & Temporal only & $\times$ \\
Closed-form posterior & \checkmark & $\times$ & $\times$ & $\times$ \\
Steps to generate (estimate) & 60--80 & 60 (per stage) & 45 & 200+ \\
\bottomrule
\end{tabularx}
\caption{Comparison of RDD-Video with existing multi-resolution video methods.}
\end{table}

\textbf{Key differentiators:}
\begin{enumerate}[nosep]
\item \textbf{Mathematical unification}: All stages are part of one continuous diffusion process; no artificial boundaries.
\item \textbf{DPM-Solver}: Can use 2nd/3rd order solvers across the entire trajectory $\to$ 60--80 total steps vs. 180+ for Pyramid-Flow.
\item \textbf{No information loss}: Drift \emph{continuously preserves} low-frequency spatiotemporal structure; renoising \emph{discards and re-injects}.
\item \textbf{Temporal coherence}: The temporal drift ensures low-frequency temporal structure (slow motion, global scene flow) is protected throughout denoising.
\end{enumerate}

%=============================================================================
\section{Experiment Design}
%=============================================================================

\subsection{Phase 1: Proof of Concept (2D$\to$Short Video)}

\begin{table}[h]
\centering
\begin{tabular}{lll}
\toprule
Experiment & Description & Metric \\
\midrule
E1.1 & RDD-Video 16 frames, 256$\times$256, UCF-101 & FVD, IS \\
E1.2 & Ablation: spatial-only vs. temporal-only vs. joint drift & FVD \\
E1.3 & Ablation: noise filling vs. renoising at transitions & FVD, temporal consistency \\
E1.4 & DPM-Solver 20/40/80 steps vs. DDPM 250 steps & FVD vs. NFE \\
\bottomrule
\end{tabular}
\end{table}

\subsection{Phase 2: Scaling Up}

\begin{table}[h]
\centering
\begin{tabular}{lll}
\toprule
Experiment & Description & Metric \\
\midrule
E2.1 & 512$\times$512, 48 frames, text-conditioned & VBench \\
E2.2 & Compare with vanilla DiT video baseline & VBench, FVD \\
E2.3 & Training efficiency: GPU-hours to reach target FVD & GPU-hours \\
E2.4 & Data-noise alignment ablation & Convergence speed \\
\bottomrule
\end{tabular}
\end{table}

\subsection{Phase 3: Full Evaluation}

\begin{table}[h]
\centering
\begin{tabular}{lll}
\toprule
Experiment & Description & Metric \\
\midrule
E3.1 & 768p, 5--10s video, comparison with Pyramid-Flow & VBench, human eval \\
E3.2 & Inference speed benchmark & Time per video, NFE \\
E3.3 & Long video generation (30s+) via sliding window & Temporal consistency \\
E3.4 & Image-to-video (I2V) extension & Quality, motion \\
\bottomrule
\end{tabular}
\end{table}

\subsection{Key Ablations}

\begin{enumerate}[nosep]
\item \textbf{Stage lengths}: [250,350,400] vs. [200,300,500] vs. [333,333,334]
\item \textbf{Temporal downsampling method}: avg-pool vs. stride vs. Gaussian blur
\item \textbf{Number of stages}: 2 vs. 3 vs. 4
\item \textbf{CFG schedule}: Constant vs. stage-dependent vs. timestep-linear
\item \textbf{Data-noise alignment}: With vs. without Hungarian matching
\end{enumerate}

%=============================================================================
\section{Implementation Roadmap}
%=============================================================================

\begin{enumerate}[nosep]
\item \textbf{Week 1--2}: Implement spatiotemporal drift in \texttt{gaussian\_diffusion\_with\_drift.py}; extend to 5D tensors $(B,C,T,H,W)$
\item \textbf{Week 2--3}: Adapt DiTMultiRes to 3D with Skiparse Attention and 3D RoPE
\item \textbf{Week 3--4}: Training pipeline: data loading (video dataset), progressive training, gradient accumulation
\item \textbf{Week 4--5}: Sampling pipeline: DPM-Solver with spatiotemporal drift correction, noise filling transitions
\item \textbf{Week 5--6}: Phase 1 experiments on UCF-101 / WebVid subset
\item \textbf{Week 6--8}: Scale up, evaluate on VBench, prepare paper figures
\end{enumerate}

%=============================================================================
\section{Conclusion}
%=============================================================================

RDD-Video extends the Resolution-Drift Diffusion framework to video generation by introducing a \textbf{spatiotemporal drift} that simultaneously encodes spatial and temporal multi-resolution information. The mathematical properties of RDD --- single Markov chain, closed-form posterior, DPM-Solver compatibility --- are fully preserved in the video setting. Compared to Pyramid-Flow and TPDiff which require explicit renoising and are limited to first-order ODE solvers, RDD-Video achieves theoretical advantages in sampling efficiency (fewer NFE), training stability (single prediction target), and generation quality (continuous information preservation without renoising loss).

The proposed approach combines the best ideas from the video generation literature: Skiparse Attention for computation efficiency (Open-Sora), progressive training (CogVideoX), data-noise alignment for faster convergence (TPDiff), and frequency-adaptive SNR (RDD's core innovation). This positions RDD-Video as a principled and efficient framework for high-quality video generation.

\end{document}
